% !TeX root = ../nc-quadric-avatar.tex
% 2026-11-22 / root-factor monodromy and generic Brauer obstruction

\section{根の内部交換と相対avatarのBrauer障害}
\label{sec:qa-root-descent-brauer}

本節では有限directed End代数の相対比較と、その上の線束を
無標識の二次形式モジュライへ降下させる問題を区別する。
根対を\emph{内部で}交換する操作には厳密な行列表示があるが、
因子の逆数倍も忘れると一般conicが非分裂になる。
とくに$E_f$と$C_q$の相違とは独立な障害が現れる。

\begin{proposition}[二つの根対の内部交換による厳密な有限End降下]
\label{prop:qa-root-swaps-linearization}
命題~\ref{prop:qa-point-relative-split-root-end}の
順序付き一次形式空間$S$上で
$W=(\mathbb Z/2)^2$を、$\ell_1,\ell_2$の交換と
$m_1,m_2$の交換で作用させる。
$W$は$S$に自由に作用する。
$\iota\in k$を$\iota^2=-1$となるように一つ選ぶと、
相対spinor加群$\mathcal X_c,\mathcal X_d$および
四項$(\mathcal B,\mathcal X_c,\mathcal B(1),
\mathcal X_d(1))$には、$W$の\emph{厳密な}
線形化を与えられる。
従って相対有限directed Endの同型
\[
 \mathfrak B_S\xrightarrow{\sim}\mathfrak G_S
\]
は$W$-同変にでき、
有限\'{e}tale商$S/W$上の有限局所自由代数の同型へ降下する。
降下後の比較も$S/W$の任意の基底変更と可換する。
根対の内部交換と、二つの二次形式の\emph{相互交換}は
異なる作用であり、後者についてはここでは述べない。
\end{proposition}
\begin{proof}
$\ell_1=s,\ell_2=t$と記し、右加群の
二項行列表示を
\[
 M(s,t)=\begin{pmatrix}c&s\\t&-c\end{pmatrix},
 \quad
 J=\begin{pmatrix}0&1\\-1&0\end{pmatrix},
 \quad J^2=-I_2
\]
とする。行列の各成分$s,t,c$は相互に可換するので
直接の乗法から
\[
 M(t,s)=JM(s,t)J,\qquad
 (\iota J)M(s,t)=M(t,s)(-\iota J).
 \tag{Swap-matrix}\label{eq:qa-root-swap-matrix}
\]
を得る。
ここで$(\iota J)^2=(-\iota J)^2=I_2$である。
従って二項表示の自由加群に対し、
余核側で$\iota J$、関係側で$-\iota J$を使えば、
二回の根交換が恒等である鎖同型が定まる。
その余核への作用が$\mathcal X_c$の厳密な
$\mathbb Z/2$-線形化である。
$\ell$の交換は$\mathcal X_d$と$\mathcal B$には恒等で作用する。
同様に$m$の交換には$d,m_1,m_2$を使い、
$\mathcal X_d$に同じ行列で線形化する。
二つの作用は異なる加群成分にのみ非自明に働くので可換し、
四項と全Hom・合成に$W$-作用を与える。

$f_1=-\ell_1\ell_2$、$f_2=-m_1m_2$は
それぞれの内部交換で不変だから、
$\mathcal B,\mathcal R=\mathcal B/(cd-dc)$には
基底変換と両立する自明な係数の輸送がある。
$\mathcal X_i\otimes_{\mathcal B}\mathcal R$の
sheafificationに上の線形化を輸送し、
その線束を$\mathcal F,\mathcal G$と同定する。
四項の後二項
$L_2\simeq\mathcal F^2$、$L_3\simeq\mathcal F^2\mathcal G$
には、それぞれ$\mathcal B(1)$と$\mathcal X_d(1)$の
制限から線形化を輸送する。
ここで$\mathcal F^2\simeq\OO_{\mathcal E}(1)$の
選択された同型が、$\mathcal F$の線形化の
tensor平方と自動的に両立するとは仮定しない。
必要なのは四つの\emph{対象ごと}に整合する線形化であり、
それにより命題~\ref{prop:qa-point-relative-split-root-end}の
有限End射は自然性から$W$-同変である。
二つの一次形式が独立なので、それぞれの交換が固定する
$S$の幾何点は存在しない。
$W$は有限定数群で標数零だから
$S\to S/W$は有限\'{e}tale torsorである。
有限局所自由加群・その積と代数同型の
faithfully flat descentが適用でき、主張が従う。

$\iota$を$-\iota$へ変えると、各交換の
線形化が符号指標だけ変わる。
従ってこの命題は線形化の\emph{存在}を述べ、
選択に依らない正準線形化とは主張しない。
\end{proof}

\begin{proposition}[普遍二次形式のconicは一般に非分裂]
\label{prop:qa-universal-conic-brauer-obstruction}
$k$を標数零の代数閉体とし、
\[
 U=\operatorname{Spec}k[a,b,d,(b^2-4ad)^{-1}],
 \quad K=k(a,b,d),\quad \Delta=b^2-4ad
\]
と置く。$K$上の滑らかなconic
\begin{equation}
 Q:\quad c^2=au^2+buv+dv^2\subset\PP^2_K
 \label{eq:qa-generic-conic}
\end{equation}
は$K$-有理点を持たない。
そのSeveri--Brauer類は
\[
 \alpha(Q)=(a,\Delta)\ne0\quad\text{in}\quad
 \operatorname{Br}(K)[2]
 \tag{Brauer-class}\label{eq:qa-conic-brauer-class}
\]
であり、$\Delta$を一様化元とする離散付値の
留数は$[a]\ne1$である。
従って$Q$には幾何次数一の線束が存在せず、
$\bar K$上の同型$Q_{\bar K}\simeq\PP^1$で定まる
次数一束$\OO_{\PP^1}(1)$は$K$上へ降下しない。
ここで射影平面から制限した$\OO_Q(1)$は
幾何次数\emph{二}である。

さらに二つの独立な一般二次形式の係数体を$K_{12}$とし、
$Q_c,Q_d$のファイバー積である滑らかな
$E=Q_c\times_{\PP^1}Q_d$を考える。
幾何学的な次数二pencil
$F_c=p_c^*\OO_{(Q_c)_{\bar K_{12}}\simeq\PP^1}(1)$と
$F_d=p_d^*\OO_{(Q_d)_{\bar K_{12}}\simeq\PP^1}(1)$は
Galois不変なPicard\emph{類}ではあるが、
一般には$E$上の実際の線束として降下しない。
そのHochschild--Serre障害はそれぞれ
\begin{equation}
 \partial_E[F_c]=\alpha(Q_c),\qquad
 \partial_E[F_d]=\alpha(Q_d),\qquad
 \partial_E[F_dF_c^{-1}]=\alpha(Q_d)-\alpha(Q_c).
 \label{eq:qa-point-pencils-brauer-obstruction}
\end{equation}
であり、三つとも$K_{12}$のBrauer群で非零である。
\end{proposition}
\begin{proof}
$K=k(a,b,\Delta)$と見なせる。
$a\ne0$の開上で
$z=u+(b/2a)v$と置くと、
\eqref{eq:qa-generic-conic}は
\[
 c^2-a z^2=-\frac{\Delta}{4a}v^2
 \tag{Diagonal-conic}\label{eq:qa-diagonal-conic}
\]
となる。右辺は$-\Delta/(4a)$を係数とする
ノルム方程式で、対応するquaternion類は
$(a,-\Delta/(4a))=(a,\Delta)$である
（$-1$と$4$は$k$で平方）。
$\Delta$の付値に対し$a$は単元で、
剰余体$k(a,b)$内の$a$は平方ではない。
従ってquaternion類の留数は非自明な$[a]$となる。

非分裂性はBrauer類の一般論を用いずに次のようにも示せる。
$\Delta$-進DVRで仮に\eqref{eq:qa-diagonal-conic}に
射影有理点があるとする。座標$(c,z,v)$をprimitiveに規格化し、
$\Delta$で還元すると
$\bar c^{\,2}=a\bar z^{\,2}$となる。
$a$が剰余体で非平方なので$\bar c=\bar z=0$である。
従って$\bar v\ne0$だが、
左辺の付値は二以上、右辺の付値は一となり矛盾する。
滑らかな種数零曲線上の幾何次数一線束が
$K$上に存在すればRiemann--Rochにより
その非零切断の零点から次数一の$K$-有効因子、
すなわち$K$-有理点が得られるので、存在しない。

二つの独立係数のときも、
第一の判別式$\Delta_c$で取った留数は
$\alpha(Q_c)$に対し$[a_c]\ne1$、
$\alpha(Q_d)$に対し零である。
同様に$\Delta_d$では役割が交換する。
よって二つのquaternion類はそれぞれ非零で、
和または差も非零である。

最後の主張は、滑らかな射影幾何的整曲線に対する
Galois descentのPicard--Brauer境界
\[
 0\longrightarrow\Pic(E)
 \longrightarrow\Pic(E_{\bar K_{12}})^{G_{K_{12}}}
 \xrightarrow{\ \partial_E\ }\operatorname{Br}(K_{12})
\]
とそのpullbackに関する自然性による。
滑らかなconicの幾何次数一Picard類の障害は
そのSeveri--Brauer類$\alpha(Q_i)$であり、
$p_i:E\to Q_i$で引き戻した類の障害も同じである。
Picard群上の差を取れば第三の式を得る。
これは$\Pic^0$の\emph{類}に相当する
Galois不変点があることと、その類の
実際の線束による代表があることの違いである。
\end{proof}

\begin{corollary}[二つの分岐根を選ぶ次数四の\'{e}tale分裂被覆]
\label{cor:qa-root-choice-etale-splitting}
$U_2$を、標数零の代数閉体上の二つの
非退化二元二次形式$(f_1,f_2)$の全パラメータ空間とする。
二つの根対の相互共有は許す。
$Z_i\subset U_2\times\PP^1$を$f_i=0$の
相対零点とし、
\[
 \widetilde U_2=Z_1\times_{U_2}Z_2
       \longrightarrow U_2
 \tag{Root-choice-cover}\label{eq:qa-root-choice-four-cover}
\]
と置く。この射は、共有根および一致根対の境界を含めて
有限\'{e}tale次数四である。
$\widetilde U_2$上の二つのconic
$Q_c,Q_d$には、それぞれ自然な分岐点section
$\sigma_c,\sigma_d$が存在し、
\[
 N_c=\OO_{Q_c}(\sigma_c),\qquad
 N_d=\OO_{Q_d}(\sigma_d)
\]
は相対次数一の実際の線束である。
従って$E=Q_c\times_{\PP^1}Q_d$上には
$F=p_c^*N_c,\ G=p_d^*N_d$という
大域的な二つの次数二pencilが存在し、
各幾何fiber上で
\[
 F^2=G^2=\OO_E(1),\qquad
 P=GF^{-1}\in\Pic^0(E)[2]\setminus\{0\}
\]
となる。これらの等式はまず各幾何fiberのPicard類の等式であり、
基底上の線束捩れを固定する正準の
平方同型を主張しない。
この被覆は一般conicのBrauer障害を明示的に消すが、
相対AS heartのpointed同値や
$\widetilde U_2\to U_2$についての
線束の降下を自動的には与えない。
\end{corollary}
\begin{proof}
各$f_i$の判別式が可逆なので
$Z_i\to U_2$は有限\'{e}tale次数二である。
二つを独立に基底変更しても有限\'{e}tale性が保たれ、
fiberの次数は$2\cdot2=4$となる。
根の共有は$f_1,f_2$の相互関係であり、
個々の$Z_i\to U_2$が\'{e}taleであることを妨げない。

選ばれた根$p_i$では$f_i(p_i)=0$なので
$Q_i:c_i^2=f_i$の上に$(p_i,c_i=0)$という
sectionが定まる。$Q_i$は相対的に滑らかなconicで、
このsectionは有効Cartier因子を与える。
従って$\OO_{Q_i}(\sigma_i)$は相対次数一。
$E$への引戻しは相対次数二である。
各幾何fiberでは$Q_i\simeq\PP^1$なので
$N_i^{\otimes2}$はconic超平面束になり、
補題~\ref{lem:qa-pointed-all-rootpair-geometry}から
$P$が非自明な多重次数零の二捩れとなる。
各式の相対的な捩れは、線束の基底からの引戻しを
標識から除く既存のavatar規約に従う。
\end{proof}

\begin{remark}[相対有限Endと無標識avatarの正確な境界]
\label{rem:qa-brauer-scope}
命題~\ref{prop:qa-root-swaps-linearization}は
\emph{根対内の有限な順序交換}を消す降下である。
二つの一次形式を
$(\ell_1,\ell_2)\mapsto(t\ell_1,t^{-1}\ell_2)$
と変える$\Gm$作用も同時に割った
二次形式の係数空間への線束の降下ではない。
実際、
\[
 D(t)=\operatorname{diag}(t,1),\qquad
 M(t\ell_1,t^{-1}\ell_2)=D(t)M(\ell_1,\ell_2)D(t)^{-1}
\]
だが、根交換の行列$J$について
\[
 J D(t)J^{-1}
   =\operatorname{diag}(1,t)
   =tD(t^{-1}) .
\]
右辺の余分なスカラーは、
線束の表示が単なる$PGL_2$作用から
自動的に$GL_2$作用へ持ち上がらないことの
有限行列的な兆候である。
実際に線束が降下しないことは
命題~\ref{prop:qa-universal-conic-brauer-obstruction}で
基礎体のBrauer類から証明した。
この障害から、\emph{別の}Morita表示、twisted sheaf、
または内在的$\mathcal H_2$の相対stack自体が
存在しないとは結論しない。
従って次の課題はBrauer類を殺す適切な
gerbe/相対AS表示の構成と、二つの根対の相互交換の
相対mutation・pointingの降下である。
\end{remark}
